% =============================================================================
% CAPÍTULO 1: APRESENTAÇÃO DO SOFTWARE
% GERADO AUTOMATICAMENTE VIA scripts/sync-md-to-latex.py A PARTIR DO MD
% =============================================================================

\aberturacapitulo{Capítulo 1}{Apresentação do software}



\subsecao{1.1 Nome e descrição do software}

\textbf{Nome do software:} CineDoom.

O CineDoom é uma plataforma web integrada, no modelo \textbf{SaaS multi-tenant}, de gerenciamento operacional, bilheteria presencial e portal do cliente, voltada para \textbf{cinemas de rua, circuitos independentes e redes de pequeno e médio porte}. O sistema permite que os cinemas informatizem seus processos com três módulos centrais:

\begin{enumerate}
\def\labelenumi{\arabic{enumi}.}
\tightlist
\item
  \textbf{Módulo do Atendente (Bilheteria Presencial):} Operação ágil de guichê físico, emissão direta de ingressos com trava atômica anti-double-booking e baixa presencial de reservas originadas na web.
\item
  \textbf{Módulo do Gerente (Back-office e Gestão):} Configuração de salas, catálogo de filmes, agendamento de sessões, gestão de operadores e acompanhamento de indicadores de ocupação.
\item
  \textbf{Portal do Cliente (Espectador Final Web):} Canal público e autenticado no qual o público consulta a programação, cadastra-se com proteção de dados LGPD, seleciona assentos no mapa interativo, solicita reservas com emissão de voucher transacional e gerencia suas reservas via web.
\end{enumerate}

A \textbf{CineSoft Solutions} é uma empresa de engenharia de software B2B especializada em soluções transacionais para o setor de entretenimento e exibição cinematográfica. O \textbf{CineDoom} posiciona-se como seu produto principal para unificar o atendimento digital (B2C) e a operação física de bilheteria (B2B) sob um único motor relacional, com isolamento multi-tenant estrito e conformidade nativa com a LGPD. O ciclo de evolução comercial do produto é organizado em \textbf{três Macroentregas} (Macroentrega 1: MVP Operacional implantado; Macroentrega 2: Pagamentos Online, Mensageria Dedicada e Escala Cloud; Macroentrega 3: Emissão Fiscal NFC-e, Bomboniere e Billing Self-Service), regido pelos processos da norma ABNT NBR ISO/IEC/IEEE 12207:2021.


\subsecao{1.2 Qual é o software? O que ele faz? Qual problema resolve?}

O CineDoom é um sistema SaaS multi-tenant que integra a operação de bilheteria, o back-office do cinema e o portal do espectador em uma única plataforma transacional.


\subsubsecao{O que ele faz?}

O CineDoom centraliza e sincroniza todas as etapas da cadeia de atendimento cinematográfico em uma interface web moderna e unificada:

\begin{itemize}
\tightlist
\item
  Mantém o catálogo de filmes em cartaz e suas especificações técnicas.
\item
  Controla a infraestrutura física de salas e layouts bidimensionais de assentos.
\item
  Agenda e disponibiliza sessões cinematográficas com prevenção automática de choques de grade.
\item
  Garante reservas online no Portal do Cliente com bloqueio atômico de assento e expiração lógica Just-in-Time (JIT).
\item
  Despacha vouchers transacionais de reserva por e-mail com instruções para pagamento presencial.
\item
  Permite o cancelamento voluntário de reservas pendentes pelo próprio cliente via web antes do limite de tolerância.
\item
  Opera a venda e a baixa presencial de ingressos na bilheteria com trava atômica ACID (\emph{SELECT FOR UPDATE}), eliminando riscos de venda dupla (\emph{double-booking}).
\item
  Apresenta métricas gerenciais consolidadas sob demanda sobre a taxa de ocupação das salas e a distribuição de assentos vendidos e reservados.
\end{itemize}


\subsubsecao{Qual problema resolve?}

Cinemas independentes frequentemente operam com sistemas legados fragmentados ou controles manuais, resultando em:

\begin{enumerate}
\def\labelenumi{\arabic{enumi}.}
\tightlist
\item
  \textbf{Conflito de assentos:} Risco crítico de venda simultânea da mesma poltrona em sessões concorridas.
\item
  \textbf{Inacessibilidade para o público:} Dificuldade do espectador em consultar a grade de horários e planejar sua ida ao cinema com antecedência.
\item
  \textbf{Assentos ociosos por no-show:} Falta de regras claras para liberação de lugares não confirmados.
\item
  \textbf{Ausência de rastreabilidade e compliance:} Descumprimento das exigências da LGPD no manuseio de dados de compradores e ausência de trilha auditável de cancelamentos.
\end{enumerate}

O CineDoom elimina esses problemas ao unificar o canal digital do cliente e a bilheteria presencial sob o mesmo motor relacional transacional e auditável.


\subsubsecao{Proposta de valor}

O CineDoom entrega um único motor transacional e auditável para os canais físico e digital. Na prática, isso reduz o risco de venda dupla, libera assentos não confirmados para reaproveitamento, diminui a dependência de controles manuais e trata os dados pessoais dos compradores conforme a LGPD.


\subsecao{1.3 Público-alvo e contexto de uso}

O público-alvo é composto por:

\begin{itemize}
\tightlist
\item
  \textbf{Público de negócio (B2B):} cinemas de rua, circuitos independentes e redes de pequeno e médio porte que precisam informatizar a bilheteria, a programação e a gestão de salas.
\item
  \textbf{Público final (B2C):} espectadores que utilizam o Portal do Cliente para consultar filmes, reservar seus assentos preferidos e gerenciar seus comprovantes antes de comparecer à sessão.
\item
  \textbf{Público operacional (usuários diretos):} gerentes e atendentes que operam o back-office e a bilheteria no dia a dia do cinema.
\item
  \textbf{Público indireto:} equipe de suporte e operação da CineSoft Solutions e provedores de infraestrutura que sustentam a plataforma.
\end{itemize}

O contexto de uso abrange:

\begin{itemize}
\tightlist
\item
  \textbf{Portal do Cliente (web desktop e mobile):} acesso pelo espectador para consulta de horários, cadastro pessoal e reserva antecipada com voucher digital.
\item
  \textbf{Ambiente de bilheteria física:} atendimento presencial ágil para venda direta ou localização e baixa de reservas prévias pelo atendente. Os indicadores de desempenho desse ambiente são definidos e verificados no Capítulo 5.
\item
  \textbf{Ambiente administrativo:} planejamento semanal de exibições, cadastro de acervo e monitoramento de lotação pela gerência.
\end{itemize}

A Figura 1 ilustra a interface do Portal do Cliente, com o catálogo de filmes em cartaz e o mapa interativo de assentos utilizado pelo espectador.

\figuracomfonte{capitulo-01/prints/print-portal-cliente.png}{0.88\textwidth}{Portal do Cliente: catálogo de filmes e mapa interativo de assentos}{CineSoft Solutions (2026)}{fig:print-portal-cliente}


\subsecao{1.4 Quem vai usar? Em que situação? Qual o ambiente de operação esperado?}


\subsubsecao{Perfis de Usuários}

\begin{itemize}
\tightlist
\item
  \textbf{Espectador / Cliente Final:} Acessa o Portal do Cliente para consultar a programação, efetuar cadastro/login com credenciais seguras, selecionar assentos no mapa interativo, reservar ingressos com expiração Just-in-Time, receber comprovante por e-mail transacional, acompanhar seu histórico pessoal em ``Minhas Reservas'' e cancelar reservas antes do prazo limite.
\item
  \textbf{Atendente:} Opera o guichê presencial. Consulta sessões, seleciona assentos, localiza reservas web por código ou CPF, confirma o recebimento presencial, emite ingressos e executa cancelamentos justificados.
\item
  \textbf{Gerente:} Administra o back-office (salas, acervo de filmes, sessões de exibição), acompanha os painéis de ocupação e gerencia os usuários com perfil ATENDENTE.
\item
  \textbf{Administrador da plataforma (perfil previsto, não disponível no MVP):} Responsável pela parametrização global, auditoria de logs e gestão de tenants. Na versão atual, essas atividades são conduzidas pela operação técnica da CineSoft Solutions; o produto entregue contempla os perfis Gerente, Atendente e Cliente.
\end{itemize}

\tabelaabnt{Matriz de Perfis de Usuários e Permissões de Acesso}{CineSoft Solutions (2026)}{tab:perfis_permissoes}{\begin{center}
\small
\begin{tabularx}{\textwidth}{l p{3.2cm} X}
\toprule
\textbf{Perfil} & \textbf{Módulo de Acesso} & \textbf{Responsabilidades e Permissões} \\
\midrule
\textbf{Cliente Final} & Portal Web Responsivo & Consulta de catálogo, seleção de poltronas, checkout e ingressos. \\
\textbf{Atendente} & Bilheteria / PDV & Venda balcão, cancelamento assistido, emissão de cortesias. \\
\textbf{Gerente} & Painel Administrativo & Gestão de salas, precificação dinâmica, relatórios de ocupação. \\
\textbf{Administrador} & Configurações Globais & Gestão de tenants SaaS, auditoria de segurança e integrações. \\
\bottomrule
\end{tabularx}
\end{center}}


\subsubsecao{Ambiente de Operação Esperado}

\begin{itemize}
\tightlist
\item
  \textbf{Hardware:} Computadores e tablets nos pontos de venda físicos; smartphones, notebooks e desktops pessoais para os espectadores no Portal do Cliente.
\item
  \textbf{Software:} Navegadores web modernos com suporte a ECMAScript atual (Google Chrome, Mozilla Firefox, Safari, Microsoft Edge).
\item
  \textbf{Conectividade:} Conexão estável com a internet, uma vez que a operação é realizada em navegador e o processamento ocorre em nuvem.
\item
  \textbf{Pagamento das Reservas:} O pagamento de ingressos reservados via web é concluído exclusivamente de forma presencial no balcão da bilheteria até 1 hora antes da sessão, eliminando a dependência de gateway financeiro online e reduzindo a superfície de fraude em transações digitais.
\end{itemize}

A Figura 2 apresenta a interface da bilheteria presencial, utilizada pelo atendente para a venda direta de ingressos e a baixa de reservas originadas no portal.

\figuracomfonte{capitulo-01/prints/print-bilheteria-atendente.png}{0.88\textwidth}{Bilheteria presencial: mapa de assentos e operação de guichê}{CineSoft Solutions (2026)}{fig:print-bilheteria-atendente}


\subsecao{1.5 Escopo e principais funcionalidades}


\subsubsecao{Escopo Central e Módulos do Sistema}

\begin{enumerate}
\def\labelenumi{\arabic{enumi}.}
\tightlist
\item
  \textbf{Módulo de Filmes:} Acervo com título, sinopse, classificação indicativa, duração e material promocional.
\item
  \textbf{Módulo de Salas e Assentos:} Mapeamento bidimensional flexível (fileiras e colunas) e controle de capacidade física.
\item
  \textbf{Módulo de Sessões:} Composição de Filme + Sala + Horário + Preço, com validação de conflitos de horário.
\item
  \textbf{Módulo de Bilheteria Presencial:} Venda atômica direta, cancelamento auditável com trilha append-only protegida por trigger PostgreSQL e baixa presencial de reservas.
\item
  \textbf{Portal do Cliente:} Consulta de sessões, autenticação com proteção de dados LGPD, mapa interativo de assentos com reserva temporária, despacho de e-mail transacional, painel de reservas e cancelamento voluntário web.
\item
  \textbf{Módulo de Ocupação:} Painel consolidado sob demanda da taxa de lotação por sessão.
\item
  \textbf{Módulo de Autenticação e Controle de Acesso:} Login com sessão JWT de 8 horas, RBAC para os perfis Gerente, Atendente e Cliente e revogação imediata do acesso de usuário inativado.
\item
  \textbf{Módulo de Gestão de Usuários:} Criação, atualização e inativação de atendentes pelo gerente, com e-mail único por tenant.
\item
  \textbf{Módulo de Conformidade e Privacidade (LGPD):} Consentimento explícito, criptografia de CPF, anonimização do titular e auditoria de cancelamentos.
\item
  \textbf{Módulo de Notificações Transacionais:} Envio de voucher de reserva por e-mail, com provedor configurável conforme o ambiente.
\end{enumerate}

A Figura 3 apresenta o painel gerencial do back-office, com a programação de sessões e os indicadores de ocupação acompanhados pela gerência.

\figuracomfonte{capitulo-01/prints/print-painel-gerencial.png}{0.88\textwidth}{Back-office gerencial: programação de sessões e indicadores de ocupação}{CineSoft Solutions (2026)}{fig:print-painel-gerencial}


\subsubsecao{Fora do escopo imediato da Macroentrega 1 (Planejados nas Macroentregas 2 e 3)}

Os seguintes módulos não compõem o escopo da \textbf{Macroentrega 1 (\texttt{Release\ 1.0\ —\ MVP\ Operacional})}, estando formalmente alocados no roadmap comercial subsequente da CineSoft Solutions:

\begin{itemize}
\tightlist
\item
  \textbf{Macroentrega 2 (\texttt{Release\ 2.0}):} Gateway de pagamento online (PIX com QR Code dinâmico e cartão de crédito via Stripe/Pagar.me), despacho de e-mails transacionais em larga escala com domínio corporativo próprio (\texttt{Resend\ Pro}) e escalonamento para banco de dados dedicado com alta disponibilidade (\texttt{Supabase/Neon\ Pro}).
\item
  \textbf{Macroentrega 3 (\texttt{Release\ 3.0}):} Emissão fiscal automatizada (NFC-e / SAT integrada à SEFAZ), gestão de bomboniere e estoque, programa de fidelidade/cupons, console administrativo global de tenants com faturamento SaaS recorrente (\emph{Self-Service Billing}) e BI financeiro avançado.
\item
  \textbf{Fora do escopo global do produto web:} Aplicativos nativos nas lojas iOS/Android (o acesso ocorre integralmente via aplicação web responsiva) e produção ou streaming de conteúdo audiovisual.
\end{itemize}


\subsecao{1.6 Restrições e premissas}


\subsubsecao{Restrições}

\begin{itemize}
\tightlist
\item
  \textbf{Pagamento Presencial Exclusivo:} O sistema não integra gateway de pagamento online (cartão/PIX web); as reservas efetuadas no portal são pagas presencialmente no guichê do cinema até 1 hora antes do início da sessão.
\item
  \textbf{Operação Conectada:} O sistema opera na nuvem sob arquitetura serverless (Vercel + Neon PostgreSQL), exigindo conectividade ativa.
\item
  \textbf{Isolamento Multi-tenant Lógico:} O compartilhamento de infraestrutura é controlado por meio do \texttt{tenant\_id}, aplicado por uma extensão de tenancy do Prisma Client com \texttt{AsyncLocalStorage}. O produto não utiliza \texttt{ROW\ LEVEL\ SECURITY} nativo do PostgreSQL.
\item
  \textbf{Sem Aplicativo Nativo:} Acesso via navegador web desktop e mobile responsivo, sem empacotamento para lojas de aplicativos.
\item
  \textbf{Conformidade LGPD:} O tratamento de dados pessoais exige consentimento explícito, criptografia de CPF (AES-256-GCM + HMAC-SHA256), minimização, anonimização do titular (Art. 18) e retenção de trilha de auditoria. O CPF não é armazenado em texto puro.
\item
  \textbf{Limites de Layout Físico:} As salas suportam até 26 fileiras (A--Z) e 50 assentos por fileira.
\item
  \textbf{Trilha de Auditoria Imutável:} Os registros de cancelamento são append-only; a proteção depende da aplicação da trigger \texttt{trg\_protect\_cancellation\_audits} no PostgreSQL.
\item
  \textbf{Dependência de Provedor de E-mail:} O envio do voucher transacional depende de um provedor configurado; em desenvolvimento, o envio é registrado no console.
\item
  \textbf{Funcionalidades Descartadas do MVP:} Streaming de vídeo, emissão fiscal externa (NFC-e), fidelidade por pontos e gestão de bomboniere, conforme detalhado na seção 1.5.
\end{itemize}


\subsubsecao{Premissas}

\begin{itemize}
\tightlist
\item
  \textbf{Concorrência de Assentos:} A trava pessimista \texttt{SELECT\ FOR\ UPDATE}, a ordenação prévia dos IDs de assento e a checagem Just-in-Time impedem, por design, conflitos de assentos entre clientes web e atendentes físicos.
\item
  \textbf{Antecedência Mínima:} As reservas online realizadas com menos de 1 hora de antecedência em relação ao início da sessão são bloqueadas pelo sistema, reservando o inventário restante para o balcão.
\item
  \textbf{Expiração Lógica:} As reservas não quitadas presencialmente até 60 minutos antes do início do filme deixam de ser confirmáveis e seus assentos passam a ser elegíveis para reaproveitamento. A expiração é lógica/Just-in-Time, sem job de varredura automática.
\item
  \textbf{Contexto de Tenant:} O tenant é resolvido a partir da sessão autenticada e todas as operações de dados são filtradas por \texttt{tenant\_id}.
\end{itemize}
